\documentclass[12pt,a4paper,oneside]{report}
\usepackage[utf8]{inputenc}
\usepackage[T1]{fontenc}
\usepackage[french]{babel}
\usepackage{mathptmx}
\usepackage{geometry,setspace,graphicx,tikz,amsmath,amssymb,xcolor}
\usetikzlibrary{arrows.meta,positioning,shapes.geometric,fit,calc,backgrounds}
\usepackage{booktabs,longtable,array,tabularx,enumitem,caption,float,fancyhdr,titlesec}
\usepackage{hyperref,listings,microtype,url}

\geometry{top=2.5cm,bottom=2.5cm,left=2.7cm,right=2.4cm}
\onehalfspacing
\setlength{\parindent}{0.7cm}
\setlength{\parskip}{0.25em}
\setlist[itemize]{itemsep=0.2em,topsep=0.4em}
\setlist[enumerate]{itemsep=0.2em,topsep=0.4em}
\definecolor{ReportNavy}{HTML}{0B233F}
\definecolor{ReportMint}{HTML}{24A89A}
\definecolor{ReportRed}{HTML}{B82025}
\definecolor{LightGray}{HTML}{F2F4F7}
\definecolor{MidGray}{HTML}{667085}
\hypersetup{
  colorlinks=true,
  linkcolor=ReportNavy,
  urlcolor=ReportMint,
  citecolor=ReportNavy,
  pdftitle={Authentification et autorisation interservices dans MindArena},
  pdfauthor={Anas Houhou et Amr Azirar}
}
\titleformat{\chapter}[display]{\bfseries\color{ReportNavy}}
  {\filleft\Large Chapitre \thechapter}{1ex}{\titlerule\vspace{1ex}\Huge}
\titleformat{name=\chapter,numberless}[display]{\bfseries\color{ReportNavy}}
  {}{0pt}{\titlerule\vspace{1ex}\Huge}
\titleformat{\section}{\Large\bfseries\color{ReportNavy}}{\thesection}{0.7em}{}
\titleformat{\subsection}{\large\bfseries\color{ReportNavy}}{\thesubsection}{0.7em}{}
\pagestyle{fancy}
\setlength{\headheight}{14pt}
\fancyhf{}
\fancyhead[L]{\small Sécurité interservices dans MindArena}
\fancyhead[R]{\small ENSIAS}
\fancyfoot[C]{\thepage}
\renewcommand{\headrulewidth}{0.4pt}
\renewcommand{\arraystretch}{1.22}
\captionsetup{font=small,labelfont=bf}
\lstset{
  basicstyle=\ttfamily\footnotesize,
  backgroundcolor=\color{LightGray},
  frame=single,
  rulecolor=\color{MidGray},
  breaklines=true,
  showstringspaces=false,
  keywordstyle=\color{ReportNavy}\bfseries,
  commentstyle=\color{MidGray},
  columns=fullflexible
}
\newcolumntype{Y}{>{\raggedright\arraybackslash}X}
\newcolumntype{P}[1]{>{\raggedright\arraybackslash}p{#1}}
\newcommand{\code}[1]{\texttt{#1}}
\newcommand{\good}{\textcolor{ReportMint}{\textbf{PASS}}}

\tikzset{
  service/.style={draw=ReportNavy,thick,rounded corners=2pt,fill=LightGray,
    minimum width=3.3cm,minimum height=1.2cm,align=center,font=\small\bfseries},
  infra/.style={draw=MidGray,rounded corners=2pt,fill=white,
    minimum width=2.8cm,minimum height=0.9cm,align=center,font=\small},
  flow/.style={-{Latex[length=2.6mm]},thick,draw=ReportMint},
  trust/.style={draw=ReportRed,dashed,rounded corners=3pt,inner sep=7pt},
  note/.style={font=\scriptsize,align=center,text=MidGray}
}

\begin{document}
\pagenumbering{gobble}

\begin{titlepage}
\begin{tikzpicture}[remember picture,overlay]
  \node[anchor=center,inner sep=0pt] at (current page.center)
    {\includegraphics[width=\paperwidth,height=\paperheight]{assets/background.png}};
\end{tikzpicture}
\vspace*{-1.45cm}
\noindent\begin{minipage}[c]{0.23\textwidth}
  \includegraphics[width=0.95\linewidth]{assets/ensias.png}
\end{minipage}%
\begin{minipage}[c]{0.48\textwidth}
  \centering
  {\normalsize\bfseries Royaume du Maroc}\\[0.12cm]
  {\small Université Mohammed V de Rabat}\\[0.12cm]
  {\small\bfseries École Nationale Supérieure d'Informatique\\et d'Analyse des Systèmes}
\end{minipage}%
\begin{minipage}[c]{0.23\textwidth}
  \raggedleft\includegraphics[width=0.92\linewidth]{assets/um5logo.png}
\end{minipage}

\vspace{2.15cm}
\begin{center}
  {\Large\bfseries Rapport de projet}\\[0.15cm]
  {\large 3\textsuperscript{e} année -- Génie Logiciel}
  \vspace{0.42cm}\rule{0.92\textwidth}{0.7pt}
  \vspace{0.42cm}
  {\LARGE\bfseries Authentification et autorisation\\[0.18cm]
  interservices dans MindArena}
  \vspace{0.42cm}\rule{0.92\textwidth}{0.7pt}
\end{center}

\vspace{0.7cm}
\begin{minipage}[t]{0.48\textwidth}
  {\large\bfseries\itshape Réalisé par :}\\[0.35cm]
  {\large Anas Houhou}\\
  {\large Amr Azirar}\\[0.8cm]
  {\large\bfseries\itshape Encadré par :}\\[0.35cm]
  {\large M.~Driss BOUZIDI}
\end{minipage}\hfill
\begin{minipage}[t]{0.43\textwidth}
  \raggedleft
  {\large\bfseries\itshape Module :}\\[0.35cm]
  {\large Ingénierie des Systèmes Sécurisés}\\[0.8cm]
  {\large\bfseries\itshape Établissement :}\\[0.35cm]
  {\large ENSIAS -- UM5 Rabat}
\end{minipage}

\vfill
\begin{center}
  \includegraphics[width=0.26\textwidth]{assets/motif-top.png}\\[-0.05cm]
  {\normalsize Période : novembre 2026}\\[0.16cm]
  {\large\bfseries Année universitaire 2026--2027}
\end{center}
\end{titlepage}

\pagenumbering{roman}

\chapter*{Remerciements}
\addcontentsline{toc}{chapter}{Remerciements}
Nous remercions l'équipe pédagogique de l'ENSIAS pour le cadre fourni par le
module d'Ingénierie des Systèmes Sécurisés et pour l'occasion d'appliquer les
principes de sécurité à une architecture distribuée concrète. Nous remercions
également toutes les personnes ayant contribué aux échanges techniques et aux
retours qui ont permis de clarifier le périmètre, les choix d'architecture et la
démarche expérimentale de ce projet.

Ce travail en binôme a été l'occasion de confronter les décisions de conception,
de relire les politiques de sécurité et de vérifier collectivement les scénarios
positifs et négatifs. Cette collaboration a joué un rôle essentiel dans la
qualité du résultat final.

\chapter*{Résumé}
\addcontentsline{toc}{chapter}{Résumé}
Les communications internes d'une architecture microservices ne doivent pas
être considérées comme fiables du seul fait qu'elles transitent sur un réseau
privé. Ce rapport présente la conception et l'implémentation d'une solution
d'authentification et d'autorisation interservices dans MindArena, une plateforme
d'apprentissage compétitif construite avec Java 21, Spring Boot et Spring Cloud.

La contribution sécurise une requête réelle faisant collaborer trois services :
Submission vérifie l'utilisateur auprès d'Identity et la politique de dépôt
auprès de Challenge avant d'enregistrer une soumission. L'authentification
repose sur des JWT de service signés avec RS256, valides pendant 120 secondes.
L'autorisation combine une identité de service enregistrée, une audience exacte
et un scope à moindre privilège. La clé privée reste exclusivement dans Identity;
les services consommateurs ne disposent que de la clé publique.

La solution est validée par des tests unitaires, une matrice de sécurité et un
scénario Docker exécuté sur des volumes neufs. Les résultats montrent le succès
du parcours autorisé, le rejet des jetons absents, malformés ou mal ciblés, la
résistance à l'usurpation d'en-têtes et l'absence d'écriture partielle lorsque
Identity ou Challenge refuse la requête.

\noindent\textbf{Mots-clés :} microservices, Zero Trust, JWT, RS256,
authentification interservices, autorisation, scopes, Spring Security.

\chapter*{Abstract}
\addcontentsline{toc}{chapter}{Abstract}
Internal communication in a microservice architecture must not be trusted merely
because it travels over a private network. This report presents the design and
implementation of service-to-service authentication and authorization in
MindArena, a competitive learning platform built with Java 21, Spring Boot and
Spring Cloud.

The contribution secures a real client request involving three services:
Submission verifies the user with Identity and checks the submission policy with
Challenge before persisting data. Authentication uses RS256-signed service JWTs
with a 120-second lifetime. Authorization combines a registered service identity,
an exact target audience and a least-privilege scope. The private key remains
exclusively in Identity while consumers receive only the public key.

Unit tests, a security matrix and a fresh Docker runtime scenario validate the
solution. Results demonstrate the successful authorized flow, rejection of
missing, malformed and incorrectly targeted tokens, protection against identity
header spoofing, and the absence of partial writes when either downstream service
refuses the request.

\noindent\textbf{Keywords:} microservices, Zero Trust, JWT, RS256,
service-to-service authentication, authorization, scopes, Spring Security.

\chapter*{Liste des acronymes}
\addcontentsline{toc}{chapter}{Liste des acronymes}
\begin{longtable}{P{2.5cm}P{11.5cm}}
\toprule
\textbf{Acronyme} & \textbf{Signification} \\
\midrule
API & Application Programming Interface \\
ABAC & Attribute-Based Access Control \\
CI & Continuous Integration \\
JWT & JSON Web Token \\
JWK(S) & JSON Web Key (Set) \\
mTLS & Mutual Transport Layer Security \\
OIDC & OpenID Connect \\
PBAC & Policy-Based Access Control \\
RBAC & Role-Based Access Control \\
ReBAC & Relationship-Based Access Control \\
RSA / RS256 & Algorithme asymétrique RSA avec SHA-256 \\
TLS & Transport Layer Security \\
\bottomrule
\end{longtable}

\tableofcontents
\listoffigures
\listoftables
\clearpage
\pagenumbering{arabic}

\chapter*{Introduction générale}
\addcontentsline{toc}{chapter}{Introduction générale}
La décomposition d'une application en microservices améliore l'autonomie des
équipes et l'indépendance des déploiements, mais multiplie les frontières de
confiance. Une requête utilisateur initialement authentifiée à l'entrée du
système déclenche souvent plusieurs appels internes. Si ces appels sont acceptés
uniquement parce qu'ils proviennent du réseau interne, la compromission d'un
service ou l'accès direct à un port suffit à contourner la politique globale.

Le principe Zero Trust propose une réponse simple à formuler mais exigeante à
mettre en œuvre : ne jamais faire confiance implicitement et toujours vérifier
l'identité et les droits associés à chaque requête. Dans le contexte d'une
architecture distribuée, l'identité de l'utilisateur ne suffit pas. Le service
appelant doit lui aussi prouver qui il est, pour quelle cible le jeton a été émis
et quelle opération précise lui est accordée.

Le présent projet applique ces principes à MindArena. Le scénario retenu est la
création d'une soumission. Pour répondre à la requête du client, Submission doit
collaborer avec Identity et Challenge. Cette chaîne répond directement à
l'exigence du sujet d'intégrer une solution d'authentification et d'autorisation
dans une application faisant intervenir trois microservices.

Notre démarche comporte quatre étapes. Nous analysons d'abord le socle existant
et les mécanismes possibles. Nous motivons ensuite le choix de JWT RS256 et d'une
autorisation par identité de service et scopes. Nous détaillons l'architecture,
l'implémentation Java et la configuration. Enfin, nous soumettons la solution à
une matrice d'essais incluant les chemins autorisés, les erreurs d'authentification,
les refus d'autorisation et l'absence d'état partiel.

\chapter{Contexte, périmètre et exigences}

\section{Présentation de MindArena}
MindArena est une plateforme d'apprentissage compétitif. Les utilisateurs
rejoignent des arènes, répondent à des challenges, soumettent des productions,
votent sur les travaux de leurs pairs, échangent dans des salons et consultent
leur classement. Le projet s'appuie sur Java 21, Spring Boot 3.3 et Spring Cloud.

L'architecture comprend notamment une API Gateway, un serveur de découverte
Eureka, un serveur de configuration, ainsi que les services Identity, Challenge,
Submission, Ranking, Notification et Chat. MariaDB assure la persistance,
Redis soutient le classement et RabbitMQ transporte les événements asynchrones.

\begin{figure}[H]
\centering
\begin{tikzpicture}[node distance=7mm and 10mm]
  \node[service] (gateway) {API Gateway\\point d'entrée};
  \node[service,below left=of gateway] (identity) {Identity\\utilisateurs et jetons};
  \node[service,below=of gateway] (submission) {Submission\\soumissions};
  \node[service,below right=of gateway] (challenge) {Challenge\\arènes et défis};
  \node[infra,below=of submission] (db) {MariaDB\\schémas par domaine};
  \node[infra,left=of db] (eureka) {Eureka\\découverte};
  \node[infra,right=of db] (mq) {RabbitMQ\\événements};
  \draw[flow] (gateway) -- (identity);
  \draw[flow] (gateway) -- (submission);
  \draw[flow] (gateway) -- (challenge);
  \draw[flow] (identity) -- (db);
  \draw[flow] (submission) -- (db);
  \draw[flow] (challenge) -- (db);
  \draw[flow,dashed] (submission) -- (identity);
  \draw[flow,dashed] (submission) -- (challenge);
  \draw[flow,dotted] (eureka) -- (submission);
  \draw[flow,dotted] (submission) -- (mq);
\end{tikzpicture}
\caption{Vue simplifiée de l'architecture MindArena}
\label{fig:mindarena-overview}
\end{figure}

\section{Limite du socle initial}
Le socle initial authentifiait l'utilisateur au niveau de la Gateway. Celle-ci
transmettait ensuite des informations d'identité aux services de domaine.
Cependant, une information placée dans un en-tête HTTP n'est pas une preuve
cryptographique lorsqu'un service est joignable directement. Une requête forgée
pouvait tenter de fournir un identifiant ou un rôle arbitraire. De plus, un
appel interne ne portait pas nécessairement une identité propre au service ni
une permission limitée à l'opération demandée.

Cette observation ne signifie pas que l'ensemble de MindArena était inutilisable.
Elle identifie une frontière de confiance pertinente pour le projet : le chemin
d'écriture d'une soumission. Nous avons choisi de sécuriser ce flux de bout en
bout plutôt que de prétendre convertir tous les appels historiques.

\section{Scénario métier retenu}
Un client authentifié envoie une requête \code{POST /api/submissions}. Submission
doit s'assurer que l'utilisateur existe toujours et que le challenge accepte les
dépôts. La réponse exige donc la collaboration de trois services de domaine :

\begin{enumerate}
  \item Submission valide le JWT utilisateur et extrait le sujet signé;
  \item Identity confirme que l'utilisateur est éligible;
  \item Challenge confirme que le challenge existe, est actif et non expiré;
  \item Submission persiste uniquement après les deux décisions positives.
\end{enumerate}

La Gateway reste le point d'entrée externe, mais n'est pas comptée comme l'un
des trois microservices métier demandés par le sujet.

\section{Exigences fonctionnelles et de sécurité}
\begin{table}[H]
\centering
\caption{Exigences du flux sécurisé}
\label{tab:requirements}
\begin{tabularx}{\textwidth}{P{1.5cm}Y Y}
\toprule
\textbf{ID} & \textbf{Exigence} & \textbf{Critère observable} \\
\midrule
E1 & Authentifier le service appelant & Sujet signé et client enregistré \\
E2 & Limiter la cible & Audience exacte pour Identity ou Challenge \\
E3 & Limiter l'opération & Scope minimal exigé à chaque endpoint \\
E4 & Séparer utilisateur et service & Validation distincte et claim \code{token\_use} \\
E5 & Échouer de façon fermée & Aucun enregistrement après un refus ou une panne \\
E6 & Réduire l'exposition & Routes internes absentes de la Gateway et politique réseau étroite \\
E7 & Tracer sans fuite & Corrélation commune, aucun bearer ni secret dans les logs \\
E8 & Reproduire les preuves & Tests automatisés et scénario Docker déterministe \\
\bottomrule
\end{tabularx}
\end{table}

\chapter{Étude des solutions et choix d'architecture}

\section{Mécanismes d'authentification interservices}

\subsection{Authentification mutuelle TLS}
mTLS authentifie les deux extrémités d'une connexion au moyen de certificats et
chiffre le transport. Son principal avantage est d'ancrer l'identité dans la
couche réseau. Il s'intègre particulièrement bien à une PKI ou à un service mesh.
En revanche, le cycle de vie des certificats, leur rotation et le diagnostic des
erreurs augmentent la charge opérationnelle. Surtout, un certificat indique qui
établit la connexion mais n'exprime pas naturellement l'opération métier précise
autorisée.

\subsection{Jetons JWT}
Un JWT signé transporte des claims vérifiables : émetteur, sujet, audience,
expiration et permissions. Avec une signature asymétrique, les récepteurs
valident localement sans détenir la capacité de créer un jeton. Cette propriété
réduit le couplage à l'émetteur et facilite la démonstration des règles. En
contrepartie, la révocation n'est pas instantanée et un jeton volé reste rejouable
jusqu'à son expiration.

\subsection{OAuth 2.0 et OpenID Connect}
OAuth 2.0 fournit un cadre standard d'autorisation et le flux client credentials
convient aux identités de service. OIDC ajoute une couche d'identité surtout
orientée utilisateur. Une solution complète telle que Keycloak ou Spring
Authorization Server apporte découverte, administration, rotation et extension,
mais son déploiement aurait déplacé le centre de ce projet vers l'exploitation
d'un produit d'identité plutôt que vers l'étude explicite des contrôles.

\begin{table}[H]
\centering
\caption{Comparaison des mécanismes d'authentification}
\label{tab:authn-comparison}
\small
\begin{tabularx}{\textwidth}{P{2.2cm}Y Y Y}
\toprule
\textbf{Mécanisme} & \textbf{Forces} & \textbf{Limites} & \textbf{Adéquation au projet} \\
\midrule
mTLS & Identité mutuelle, chiffrement, contrôle au transport & PKI et rotation plus lourdes; permission métier non exprimée seule & Excellent complément de production \\
JWT RS256 & Validation locale; claims visibles; audience et scopes; clé privée non distribuée & Révocation différée; protection du bearer nécessaire & Choisi pour sa lisibilité et sa testabilité \\
OAuth 2.0/OIDC & Standard, écosystème riche, flux client credentials & Infrastructure et configuration supplémentaires & Évolution naturelle pour industrialiser \\
\bottomrule
\end{tabularx}
\end{table}

\section{Modèles de contrôle d'accès}
RBAC attribue les permissions à des rôles. ABAC prend en compte des attributs de
l'appelant, de la ressource et du contexte. PBAC formalise la décision dans des
politiques, souvent évaluées par un moteur dédié. ReBAC raisonne sur les relations
entre entités. Le scénario étudié comporte un appelant stable et deux opérations
précises; un graphe de relations ou un moteur de règles complet serait donc
disproportionné.

\begin{table}[H]
\centering
\caption{Comparaison des modèles d'autorisation}
\label{tab:authz-comparison}
\small
\begin{tabularx}{\textwidth}{P{1.8cm}Y Y}
\toprule
\textbf{Modèle} & \textbf{Atout principal} & \textbf{Appréciation pour MindArena} \\
\midrule
RBAC & Simple, stable et auditable & Utile pour l'identité de service, mais trop large seul \\
ABAC & Décision contextuelle très fine & Puissant mais complexe pour deux opérations fixes \\
PBAC & Gouvernance centralisée des règles & Pertinent à grande échelle, surdimensionné ici \\
ReBAC & Adapté aux graphes de partage & Peu pertinent pour une relation appelant-cible statique \\
Identité + scopes & Rôle stable et permission précise & Choix retenu : lisible, minimal et testable \\
\bottomrule
\end{tabularx}
\end{table}

\section{Décision retenue}
L'authentification repose sur des JWT de service signés avec RS256. Identity est
l'unique émetteur et conserve la clé privée. Submission et Challenge reçoivent
uniquement la clé publique. La durée de vie par défaut est de 120 secondes.

L'autorisation combine l'identité de service contenue dans le sujet, l'audience
cible et un scope d'opération. Cette approche peut être vue comme un RBAC de
service affiné par des permissions. Toute combinaison non déclarée est refusée.
Nous retenons mTLS comme complément de transport recommandé, non comme substitut
à l'autorisation applicative.

\section{Critères de décision}
Le choix a été évalué selon la robustesse, la lisibilité pédagogique, la facilité
d'audit, le coût d'exploitation local, la performance, la capacité de tests
négatifs et la trajectoire d'évolution. RS256 fournit le meilleur compromis pour
une réalisation en binôme et pour un rapport où chaque hypothèse doit être
observable dans le code et dans les résultats.

\chapter{Conception de la solution}

\section{Répartition des responsabilités}
\begin{table}[H]
\centering
\caption{Responsabilités et secrets par composant}
\label{tab:responsibilities}
\small
\begin{tabularx}{\textwidth}{P{2.5cm}Y Y}
\toprule
\textbf{Composant} & \textbf{Responsabilité} & \textbf{Matériel sensible} \\
\midrule
Identity & Enregistre le client, vérifie ses credentials, émet les JWT et protège la vérification utilisateur & Clé privée RSA et digest SHA-256 du secret client \\
Submission & Valide le JWT utilisateur, obtient un jeton par audience et orchestre les décisions avant écriture & Secret client brut et clé publique RSA \\
Challenge & Valide le JWT de service et la politique du challenge & Clé publique RSA seulement \\
Gateway & Valide le client externe, assainit la corrélation et ne route pas les endpoints internes & Clé publique RSA seulement \\
\bottomrule
\end{tabularx}
\end{table}

\begin{figure}[H]
\centering
\resizebox{0.97\textwidth}{!}{%
\begin{tikzpicture}[node distance=9mm and 13mm]
  \node[service] (client) {Client\\JWT utilisateur};
  \node[service,right=of client] (gateway) {API Gateway\\validation externe};
  \node[service,right=of gateway] (submission) {Submission\\orchestrateur};
  \node[service,above right=of submission] (identity) {Identity\\émetteur + décision};
  \node[service,below right=of submission] (challenge) {Challenge\\décision métier};
  \draw[flow] (client) -- node[above,note]{requête} (gateway);
  \draw[flow] (gateway) -- node[above,note]{JWT utilisateur} (submission);
  \draw[flow] (submission) -- node[above left,note]{JWT service\\audience Identity} (identity);
  \draw[flow] (submission) -- node[below left,note]{JWT service\\audience Challenge} (challenge);
  \begin{scope}[on background layer]
    \node[trust,fit=(identity)(challenge),label={[ReportRed]right:zone interne protégée}] {};
  \end{scope}
\end{tikzpicture}
}
\caption{Frontières de confiance du flux sécurisé}
\label{fig:trust-boundaries}
\end{figure}

\section{Contrat du jeton de service}
Le jeton contient les claims suivants :
\begin{itemize}
  \item \code{iss=mindarena-identity-service} : émetteur attendu;
  \item \code{sub=mindarena-submission-service} : identité de l'appelant;
  \item \code{aud} : service destinataire exact;
  \item \code{scope} : opération permise;
  \item \code{token\_use=service} : séparation avec les jetons utilisateur;
  \item \code{iat} et \code{exp} : fenêtre temporelle courte.
\end{itemize}

\begin{lstlisting}[caption={Exemple conceptuel de claims pour Challenge},label={lst:claims}]
{
  "iss": "mindarena-identity-service",
  "sub": "mindarena-submission-service",
  "aud": ["mindarena-challenge-service"],
  "scope": "challenge:submission:validate",
  "token_use": "service",
  "iat": 1795450000,
  "exp": 1795450120
}
\end{lstlisting}

\section{Politique d'autorisation}
\begin{table}[H]
\centering
\caption{Matrice d'autorisation à moindre privilège}
\label{tab:policy}
\begin{tabularx}{\textwidth}{P{3.1cm}P{2.2cm}P{4.2cm}Y}
\toprule
\textbf{Appelant} & \textbf{Cible} & \textbf{Audience} & \textbf{Scope exigé} \\
\midrule
mindarena-submission-service & Identity & mindarena-identity-service & identity:user:verify \\
mindarena-submission-service & Challenge & mindarena-challenge-service & challenge:submission:validate \\
\bottomrule
\end{tabularx}
\end{table}

Un endpoint interne exige simultanément un jeton correctement signé, le bon
émetteur, le type service, le sujet Submission, sa propre audience et le scope
exact. Un jeton destiné à Identity ne peut donc pas être réutilisé sur Challenge.

\section{Séquence complète}
\begin{figure}[H]
\centering
\begin{tikzpicture}[x=2.65cm,y=0.77cm,font=\scriptsize]
  \foreach \name/\x in {Client/0,Gateway/1,Submission/2,Identity/3,Challenge/4}
    {\node[font=\small\bfseries,text=ReportNavy] (\name) at (\x,0) {\name};
     \draw[MidGray] (\x,-0.3) -- (\x,-9.8);}
  \draw[flow] (0,-1) -- node[above]{POST + JWT utilisateur} (1,-1);
  \draw[flow] (1,-2) -- node[above]{JWT propagé} (2,-2);
  \draw[flow] (2,-3) -- node[above]{obtenir token Identity} (3,-3);
  \draw[flow] (3,-4) -- node[above]{token aud=Identity} (2,-4);
  \draw[flow] (2,-5) -- node[above]{vérifier utilisateur} (3,-5);
  \draw[flow] (3,-6) -- node[above]{éligible} (2,-6);
  \draw[flow] (2,-7) -- node[above]{valider politique + token ciblé} (4,-7);
  \draw[flow] (4,-8) -- node[above]{challenge ouvert} (2,-8);
  \node[note,text=ReportNavy] at (2,-8.75) {persistance transactionnelle};
  \draw[flow] (2,-9.5) -- node[above]{200 + corrélation} (1,-9.5);
  \draw[flow] (1,-10.2) -- node[above]{réponse} (0,-10.2);
\end{tikzpicture}
\caption{Séquence de création d'une soumission autorisée}
\label{fig:sequence}
\end{figure}

\section{Atomicité et comportement d'erreur}
Les vérifications distantes précèdent toute création de soumission. Le service
met à jour ses projections locales et persiste dans une même transaction
seulement après les deux accords. Un refus Identity ou Challenge ne laisse donc
aucune ligne partielle.

\begin{table}[H]
\centering
\caption{Sémantique des erreurs}
\label{tab:error-semantics}
\begin{tabularx}{\textwidth}{P{2cm}Y Y}
\toprule
\textbf{HTTP} & \textbf{Situation} & \textbf{Effet} \\
\midrule
401 & Jeton absent, malformé, expiré, forgé ou mauvais type & Aucune logique métier \\
403 & Jeton valide mais sujet, audience ou scope incorrect & Autorisation refusée \\
404/422 & Utilisateur ou challenge non éligible & Refus métier, aucune soumission \\
502 & Autorité de service configurée mais rejetée & Erreur de configuration interservices \\
503 & Émetteur ou service de validation indisponible & Échec fermé, aucune écriture \\
\bottomrule
\end{tabularx}
\end{table}

\chapter{Implémentation Java et configuration}

\section{Émission dans Identity}
Identity expose un endpoint interne d'émission. Submission s'authentifie avec
un identifiant et un secret client généré aléatoirement. Identity ne stocke pas
le secret brut : il calcule SHA-256 de la valeur reçue et compare le digest au
digest configuré en temps constant. Le registre détermine ensuite si le client
est autorisé à demander l'audience cible et quel scope doit être placé dans le
jeton.

\begin{itemize}
  \item \code{JwtService} construit les claims et signe avec la clé privée RSA;
  \item \texttt{ServiceClient}\allowbreak\texttt{Registry} porte la relation
    client-audience-scope;
  \item \texttt{ServiceToken}\allowbreak\texttt{Controller} traite
    l'authentification du client et renvoie le bearer à courte durée de vie.
\end{itemize}

\section{Validation dans Identity et Challenge}
Les endpoints \code{/internal/**} utilisent une chaîne Spring Security Resource
Server distincte. Le décodeur vérifie la signature, l'émetteur, l'expiration et
\code{token\_use=service}. Un \code{AuthorizationManager} vérifie ensuite le sujet,
l'audience et le scope attendus par l'endpoint.

Ce découpage distingue correctement l'authentification de l'autorisation. Une
signature invalide produit 401; une identité authentique mais insuffisamment
autorisée produit 403.

\section{Orchestration dans Submission}
\code{ServiceTokenProvider} demande un token séparé pour chaque audience et le
conserve jusqu'à quinze secondes avant son expiration. Les caches ne sont jamais
partagés entre Identity et Challenge. \code{IdentityServiceClient} et
\code{ChallengeServiceClient} ajoutent le bearer ciblé et propagent l'identifiant
de corrélation.

Avant ce flux, Submission agit également comme Resource Server pour le JWT
utilisateur. \code{UserTokenValidator} exige l'émetteur, l'audience
\code{mindarena-api} et \code{token\_use=user}. Le contrôleur extrait l'identifiant
du sujet signé; les en-têtes \code{X-User-*} ne participent pas à l'autorisation.

\section{Observabilité}
La Gateway crée ou assainit \code{X-Correlation-Id}. Le même identifiant est
propagé par Submission et placé dans le contexte de logs des trois services.
Cette corrélation permet d'expliquer une décision distribuée sans écrire les
bearer tokens, credentials Basic ou hashes de secrets dans les journaux.

\section{Configuration locale}
Les secrets sont externalisés dans un fichier \code{.env} ignoré par Git. Le
dépôt fournit seulement \code{.env.example} avec des valeurs factices. Deux
scripts génèrent les clés et credentials sans afficher la clé privée dans le
terminal.

\begin{lstlisting}[caption={Initialisation locale sécurisée},label={lst:local-setup}]
cp .env.example .env
# Remplacer les mots de passe DB, RabbitMQ et Grafana
bash scripts/generate-jwt-keys.sh --write-env .env
bash scripts/generate-service-credentials.sh --write-env .env
docker compose up -d --build
\end{lstlisting}

\begin{table}[H]
\centering
\caption{Variables principales de sécurité}
\label{tab:variables}
\small
\begin{tabularx}{\textwidth}{P{5.0cm}P{2.7cm}Y}
\toprule
\textbf{Variable} & \textbf{Consommateur} & \textbf{Rôle} \\
\midrule
\code{JWT\_PRIVATE\_KEY} & Identity & Signature RS256 uniquement \\
\code{JWT\_PUBLIC\_KEY} & Gateway, Submission, Challenge & Validation sans capacité d'émission \\
\nolinkurl{SERVICE_JWT_EXPIRATION_SECONDS} & Identity & Durée de vie, 120 s par défaut \\
\nolinkurl{SUBMISSION_SERVICE_CLIENT_ID} & Identity, Submission & Identité stable du service \\
\nolinkurl{SUBMISSION_SERVICE_CLIENT_SECRET} & Submission & Credential brut au runtime \\
\nolinkurl{SUBMISSION_SERVICE_CLIENT_SECRET_SHA256} & Identity & Digest du credential \\
\bottomrule
\end{tabularx}
\end{table}

\section{Défense réseau et déploiement}
Les routes internes ne sont pas déclarées dans la Gateway. En Kubernetes, une
NetworkPolicy autorise précisément Submission à atteindre Identity et Challenge.
Le réseau réduit l'exposition, tandis que le JWT maintient l'autorisation même si
un acteur parvient à joindre l'endpoint.

\begin{figure}[H]
\centering
\resizebox{0.94\textwidth}{!}{%
\begin{tikzpicture}[node distance=9mm and 9mm]
  \node[service] (gw) {Gateway\\entrée publique};
  \node[service,below left=of gw] (sub) {Submission\\clé publique + secret client};
  \node[service,below right=of gw] (id) {Identity\\clé privée + digest};
  \node[service,right=of id] (ch) {Challenge\\clé publique};
  \node[infra,below=of id] (secret) {Secrets Kubernetes / .env local};
  \draw[flow] (gw) -- (sub);
  \draw[flow] (sub) -- node[above,note]{autorisé} (id);
  \draw[flow] (sub) to[bend right=18] node[below,note]{autorisé} (ch);
  \draw[flow,dashed] (secret) -- (sub);
  \draw[flow,dashed] (secret) -- (id);
  \draw[ReportRed,very thick,-{Latex}] (gw) to[bend left=25] node[above,note,text=ReportRed]{aucune route /internal} (id);
\end{tikzpicture}
}
\caption{Déploiement et placement des secrets}
\label{fig:deployment}
\end{figure}

\chapter{Validation et résultats}

\section{Stratégie de validation}
La validation combine plusieurs niveaux. Les tests unitaires isolent les
validateurs de claims, le registre client, la politique d'autorisation, la règle
de challenge et l'orchestration transactionnelle. Les builds Maven vérifient
l'intégration de chaque service. Enfin, le script runtime effectue de vrais
appels HTTP avec de vrais JWT sur une pile Docker et interroge MariaDB pour
prouver l'absence d'écriture partielle.

\section{Matrice d'essais automatisée}
\begin{longtable}{P{7.4cm}P{2.2cm}P{2.2cm}P{1.8cm}}
\caption{Résultats de la matrice runtime}\label{tab:runtime-results}\\
\toprule
\textbf{Scénario} & \textbf{Cible} & \textbf{Attendu} & \textbf{Résultat} \\
\midrule
\endfirsthead
\toprule
\textbf{Scénario} & \textbf{Cible} & \textbf{Attendu} & \textbf{Résultat} \\
\midrule
\endhead
Aucun token & Identity interne & 401 & \good \\
Bearer malformé & Identity interne & 401 & \good \\
JWT utilisateur sur endpoint service & Identity interne & 401 & \good \\
En-têtes \code{X-User-*} usurpés & Identity interne & 401 & \good \\
En-têtes usurpés sur appel direct & Submission & 401 & \good \\
JWT service sur endpoint utilisateur & Submission & 401 & \good \\
Token audience Identity réutilisé & Challenge & 403 & \good \\
Token Identity correct & Identity & 200 & \good \\
Token Challenge correct & Challenge & 200 & \good \\
Endpoint interne via Gateway & Gateway & 404 & \good \\
Identity refuse un utilisateur supprimé & Flux complet & 422 & \good \\
Challenge refuse un délai expiré & Flux complet & 422 & \good \\
Nombre de soumissions après refus & MariaDB & inchangé & \good \\
Utilisateur valide et challenge ouvert & Flux complet & 200 & \good \\
\bottomrule
\end{longtable}

\section{Commandes reproductibles}
L'override de test publie Identity, Challenge et Submission uniquement sur
l'interface loopback. Les volumes du projet isolé sont supprimés après l'essai.

\begin{lstlisting}[caption={Exécution du scénario complet},label={lst:runtime}]
docker compose -p mindarena-security-qa \
  -f docker-compose.yml \
  -f docker-compose.security-test.yml \
  up -d --build

COMPOSE_PROJECT_NAME=mindarena-security-qa \
  bash scripts/smoke-service-security.sh

docker compose -p mindarena-security-qa \
  -f docker-compose.yml \
  -f docker-compose.security-test.yml \
  down -v
\end{lstlisting}

La dernière ligne observée est :
\begin{lstlisting}
Runtime service-security smoke checks passed.
\end{lstlisting}

\section{Validation de la livraison}
Le code a été archivé à partir des fichiers suivis par Git, sans \code{.env},
\code{target/}, \code{.idea}, clés privées ni métadonnées Git. L'archive a été
extraite dans un répertoire temporaire neuf. Depuis cette extraction, la matrice
de sécurité, les neuf builds de microservices et la vérification Maven racine ont
réussi. Le projet racine a exécuté 11 tests sans échec.

Cette étape apporte une preuve plus forte qu'un build dans l'espace de travail :
elle vérifie que l'archive livrée contient effectivement toutes les sources et
configurations nécessaires.

\section{Traçabilité des exigences}
\begin{table}[H]
\centering
\caption{Couverture des exigences}
\label{tab:traceability}
\small
\begin{tabularx}{\textwidth}{P{1.4cm}Y Y}
\toprule
\textbf{ID} & \textbf{Contrôle implémenté} & \textbf{Preuve} \\
\midrule
E1 & Client enregistré, sujet signé, secret vérifié & Tests registre + jeton valide \\
E2 & Token distinct par audience & Réutilisation Identity sur Challenge refusée 403 \\
E3 & Scope exact par endpoint & Tests des AuthorizationManager \\
E4 & \code{token\_use} et validateurs distincts & User sur service et service sur user refusés 401 \\
E5 & Vérifications avant transaction & Deux refus, nombre de lignes inchangé \\
E6 & Aucune route interne + NetworkPolicy & Gateway 404 et manifeste Kubernetes \\
E7 & Corrélation propagée, aucun bearer loggé & Logs corrélés vérifiés sur le scénario \\
E8 & Scripts CI et smoke déterministes & Builds et extraction propre réussis \\
\bottomrule
\end{tabularx}
\end{table}

\chapter{Analyse de sécurité, limites et perspectives}

\section{Modèle de menaces}
\begin{longtable}{P{3.2cm}P{6.1cm}P{4.2cm}}
\caption{Menaces, contrôles et risques résiduels}\label{tab:threat-model}\\
\toprule
\textbf{Menace} & \textbf{Contrôle} & \textbf{Risque résiduel} \\
\midrule
\endfirsthead
\toprule
\textbf{Menace} & \textbf{Contrôle} & \textbf{Risque résiduel} \\
\midrule
\endhead
Usurpation de service & Client enregistré, secret entropique, digest SHA-256 et comparaison en temps constant & Rotation après suspicion de fuite \\
Vol ou rejeu de jeton & Durée 120 s, audience étroite, scope minimal, aucun token loggé & Rejeu possible durant la courte fenêtre \\
Élévation de privilège & Rejet par défaut et combinaison exacte sujet-audience-scope & Revue nécessaire à chaque nouveau service \\
Accès direct interne & Resource Server sur chaque endpoint et NetworkPolicy & L'accès hôte/cluster reste un sujet infrastructure \\
Confusion utilisateur/service & Claim \code{token\_use}, audience et validateurs distincts & Clés séparées possibles en durcissement futur \\
Usurpation d'en-tête & Sujet du JWT utilisateur, jamais \code{X-User-*} & Ne pas publier les ports métier \\
Exposition de clé & Clé privée injectée uniquement dans Identity & HSM ou secret manager recommandé \\
Panne aval & 503 cohérent et aucune écriture avant accord & Circuit breaker à ajouter à grande échelle \\
\bottomrule
\end{longtable}

\section{Avantages de la solution}
\begin{itemize}
  \item la validation locale évite un appel d'introspection à chaque requête;
  \item les permissions sont visibles dans le jeton, le code et les tests;
  \item la clé privée n'est jamais distribuée aux services consommateurs;
  \item l'audience empêche le rejeu transversal entre services;
  \item les erreurs 401, 403, 422, 502 et 503 restent sémantiquement distinctes;
  \item la solution s'intègre aux primitives standard de Spring Security.
\end{itemize}

\section{Limites actuelles}
Un JWT autoporté n'offre pas une révocation instantanée. La courte durée de vie
réduit le risque mais ne l'annule pas. La clé publique et les credentials sont
injectés statiquement; la rotation demande une opération coordonnée. Le transport
local n'est pas présenté comme un déploiement de production et doit être protégé
par TLS. Enfin, le périmètre validé est le chemin d'écriture Submission; les
autres appels historiques devront adopter progressivement le même modèle.

\section{Perspectives d'industrialisation}
\begin{enumerate}
  \item activer TLS et éventuellement mTLS entre tous les services;
  \item publier un JWKS, utiliser \code{kid} et automatiser la rotation avec
    une période de chevauchement;
  \item déplacer les secrets vers Vault ou un gestionnaire cloud;
  \item adopter Spring Authorization Server ou Keycloak pour un OAuth 2.0 complet;
  \item ajouter timeouts, circuit breakers, métriques et alertes d'autorisation;
  \item gouverner un catalogue de scopes et automatiser les revues de politique;
  \item étendre le modèle aux autres flux de MindArena sans élargir
    implicitement les droits existants.
\end{enumerate}

\chapter*{Conclusion générale}
\addcontentsline{toc}{chapter}{Conclusion générale}
Ce projet a transformé une faiblesse concrète de confiance implicite en une
chaîne d'autorisation interservices explicite et testable. Pour créer une
soumission, Submission doit désormais prouver son identité, présenter un jeton
destiné à chaque service cible et disposer du scope exact. Identity et Challenge
valident indépendamment la preuve cryptographique et la permission avant de
prendre leur décision métier.

Le choix de JWT RS256 répond au besoin de validation locale tout en conservant
la capacité de signature dans Identity. L'association identité-audience-scope
met en œuvre le moindre privilège sans introduire un moteur de politique
disproportionné. La séparation des jetons utilisateur et service, le rejet des
en-têtes usurpés, l'absence de routes internes publiques et la règle de non-
persistance avant les deux accords matérialisent le principe Zero Trust.

La validation ne se limite pas au chemin heureux. Les cas d'attaque et de refus
sont automatisés, le nombre de lignes est contrôlé après les échecs et l'archive
de livraison a été reconstruite puis testée depuis une extraction propre. Ces
preuves permettent de distinguer clairement l'implémentation, la vérification
automatique et les améliorations proposées pour une exploitation industrielle.

\begin{thebibliography}{9}
\bibitem{rfc7519}
M. Jones, J. Bradley et N. Sakimura,
\emph{JSON Web Token (JWT)}, RFC 7519, IETF, 2015.
\url{https://www.rfc-editor.org/rfc/rfc7519}

\bibitem{rfc6749}
D. Hardt,
\emph{The OAuth 2.0 Authorization Framework}, RFC 6749, IETF, 2012.
\url{https://www.rfc-editor.org/rfc/rfc6749}

\bibitem{rfc8705}
B. Campbell et al.,
\emph{OAuth 2.0 Mutual-TLS Client Authentication and Certificate-Bound Access Tokens},
RFC 8705, IETF, 2020.
\url{https://www.rfc-editor.org/rfc/rfc8705}

\bibitem{nist-zero-trust}
S. Rose et al.,
\emph{Zero Trust Architecture}, NIST Special Publication 800-207, 2020.
\url{https://csrc.nist.gov/publications/detail/sp/800-207/final}

\bibitem{spring-security}
VMware,
\emph{Spring Security Reference -- OAuth 2.0 Resource Server JWT}.
\url{https://docs.spring.io/spring-security/reference/servlet/oauth2/resource-server/jwt.html}

\bibitem{owasp-microservices}
OWASP Foundation,
\emph{Microservices Security Cheat Sheet}.
\url{https://cheatsheetseries.owasp.org/cheatsheets/Microservices_Security_Cheat_Sheet.html}
\end{thebibliography}

\appendix
\chapter{Correspondance avec le sujet}
\begin{longtable}{P{5.1cm}P{6cm}P{2.6cm}}
\toprule
\textbf{Demande} & \textbf{Réponse MindArena} & \textbf{Emplacement} \\
\midrule
Solution d'authentification/autorisation & JWT RS256, identité de service, audience et scopes & Chapitres 2 à 4 \\
Trois microservices pour une requête & Submission, Identity et Challenge & Chapitres 1 et 3 \\
Avantages, inconvénients et critères & Comparaison JWT, mTLS, OAuth2, RBAC, ABAC, PBAC et ReBAC & Chapitres 2 et 6 \\
Architecture et fonctionnement & Placement des clés, claims, politique, séquence et erreurs & Chapitre 3 \\
Installation et configuration & Java 21, IntelliJ, secrets, Docker Compose & Chapitre 4 et annexe B \\
Implémentation Java importable & Spring Boot, POM Maven par service, archive contrôlée & Chapitres 4 et 5 \\
\bottomrule
\end{longtable}

\chapter{Guide court de reproduction}
\section{Import IntelliJ}
Extraire l'archive, ouvrir le répertoire \code{MindArena}, sélectionner Java 21
comme SDK et laisser IntelliJ importer le \code{pom.xml} racine ainsi que les
projets Maven situés sous \nolinkurl{services/*/pom.xml}. Aucun fichier \code{.idea}
ni chemin spécifique à une machine n'est nécessaire.

\section{Préparation et tests}
\begin{lstlisting}
cp .env.example .env
# Renseigner les mots de passe locaux
bash scripts/generate-jwt-keys.sh --write-env .env
bash scripts/generate-service-credentials.sh --write-env .env

bash scripts/test-service-security.sh
bash scripts/build-microservices.sh
mvn -B verify
\end{lstlisting}

\section{Fichiers de référence}
\begin{itemize}
  \item \code{docs/security/EVALUATOR\_QUICKSTART.md} : parcours évaluateur;
  \item \code{docs/security/service-to-service-security.md} : conception technique;
  \item \code{docs/report/VALIDATION\_EVIDENCE.md} : résultats détaillés;
  \item \code{docs/report/PRESENTATION\_GUIDE.md} : démonstration et questions;
  \item \code{scripts/smoke-service-security.sh} : preuve runtime;
  \item \code{scripts/create-delivery-archive.sh} : création de l'archive propre.
\end{itemize}

\chapter{Checklist de remise}
\begin{itemize}[label=$\square$]
  \item vérifier les noms du binôme et de l'encadrant sur la couverture;
  \item ouvrir le PDF final et contrôler les diagrammes, tableaux et références;
  \item importer l'archive dans IntelliJ avec Java 21;
  \item confirmer l'absence de \code{.env}, clés privées, \code{target/} et \code{.idea};
  \item nommer clairement le PDF et l'archive avec les deux membres du binôme;
  \item déposer les deux fichiers sur Moodle avant le 24 novembre 2026 à minuit;
  \item télécharger ou prévisualiser les fichiers déposés et conserver la confirmation.
\end{itemize}

\end{document}
